\subsubsection{口交中的射精、吞咽与精液安全}

口交要不要"含到最后"、能不能吞咽，是咨询中出现频率很高却少有人公开讨论的问题。这里给出基于医学事实的说明，帮助双方做出\textbf{知情且自愿}的选择。

\vspace{0.5em}
\noindent\textbf{先讲清基本事实}：

\begin{itemize}
	\item \textbf{精液的成分}：精子只占很小一部分，其余主要是精浆——水、果糖、蛋白质、酶、锌、维生素 C 等。对\textbf{健康}的伴侣而言，吞咽精液\textbf{本身没有危害}，胃酸会消化其中的成分，"吞精伤身""一次吞咽等于吸收多少营养"之类的说法都缺乏依据
	\item \textbf{但"入口"就意味着风险敞口}：口交是低风险但\textbf{非零风险}的性传播感染途径。淋病、衣原体、梅毒、疱疹与 HPV 都可经口--生殖器接触传播；若对方是 HIV 感染者且未接受有效抗病毒治疗，口内射精（尤其是射在口腔黏膜破损处）的风险高于"只含不射"
	\item \textbf{射在体外不等于绝对安全}：前列腺液与尿道球腺液在射精前就可能携带病原体，"没射就没事"是常见误区
\end{itemize}

\vspace{0.5em}
\noindent\textbf{降低风险的实用做法}：

\begin{itemize}
	\item \textbf{知道对方健康状况再决定}：固定且双方检测过的伴侣之间，是否吞咽属于纯粹的个人偏好；对健康状况不明的伴侣，应坚持使用避孕套口交——这也是"口交用套"最实际的理由之一
	\item \textbf{注意口腔与生殖器的完整}：口腔溃疡、牙龈出血、咽喉炎症，或生殖器有伤口与异常分泌物时，应暂停此类活动（详见上一小节的安全与卫生清单）
	\item \textbf{射精前给出明确信号}：无论选择体外、口外还是其他方式，提前沟通、把选择权交给双方，是口交中"射精礼仪"的核心
\end{itemize}

\vspace{0.5em}
\noindent\textbf{两个常见疑问}：

\begin{itemize}
	\item \textbf{"吞咽后担心怀孕"}：经消化道不会怀孕——精子无法穿过胃肠到达子宫。但口交后继续阴道接触时，手上或生殖器上沾到的精液\textbf{仍可能}导致怀孕与感染，需要清洁或更换避孕套
	\item \textbf{"闻到/尝到味道很冲"}：精液气味受饮食（芦笋、大蒜、烟酒）、饮水与禁欲天数影响，属正常现象；若伴有恶臭、颜色异常（黄绿、带血），提示生殖道感染，应先就医再继续无保护的口交
\end{itemize}

\vspace{0.5em}
\noindent\textbf{关于"意愿"的一句话}：吞不吞、吐不吐、射在哪里，都\textbf{没有标准答案}——愿意是情趣，不愿意是权利。任何一方被催促、施压后的"勉强答应"，都会在事后留下不舒服的余味；把偏好尽早摆到桌面上谈，比事后别扭好得多。

